% Options for packages loaded elsewhere
\PassOptionsToPackage{unicode}{hyperref}
\PassOptionsToPackage{hyphens}{url}
\documentclass[
  a4paper]{article}
\usepackage{xcolor}
\usepackage[margin=25mm]{geometry}
\usepackage{amsmath,amssymb}
\setcounter{secnumdepth}{-\maxdimen} % remove section numbering
\usepackage{iftex}
\ifPDFTeX
  \usepackage[T1]{fontenc}
  \usepackage[utf8]{inputenc}
  \usepackage{textcomp} % provide euro and other symbols
\else % if luatex or xetex
  \usepackage{unicode-math} % this also loads fontspec
  \defaultfontfeatures{Scale=MatchLowercase}
  \defaultfontfeatures[\rmfamily]{Ligatures=TeX,Scale=1}
\fi
\usepackage{lmodern}
\ifPDFTeX\else
  % xetex/luatex font selection
\fi
% Use upquote if available, for straight quotes in verbatim environments
\IfFileExists{upquote.sty}{\usepackage{upquote}}{}
\IfFileExists{microtype.sty}{% use microtype if available
  \usepackage[]{microtype}
  \UseMicrotypeSet[protrusion]{basicmath} % disable protrusion for tt fonts
}{}
\makeatletter
\@ifundefined{KOMAClassName}{% if non-KOMA class
  \IfFileExists{parskip.sty}{%
    \usepackage{parskip}
  }{% else
    \setlength{\parindent}{0pt}
    \setlength{\parskip}{6pt plus 2pt minus 1pt}}
}{% if KOMA class
  \KOMAoptions{parskip=half}}
\makeatother
\usepackage{color}
\usepackage{fancyvrb}
\newcommand{\VerbBar}{|}
\newcommand{\VERB}{\Verb[commandchars=\\\{\}]}
\DefineVerbatimEnvironment{Highlighting}{Verbatim}{commandchars=\\\{\}}
% Add ',fontsize=\small' for more characters per line
\newenvironment{Shaded}{}{}
\newcommand{\AlertTok}[1]{\textcolor[rgb]{1.00,0.00,0.00}{\textbf{#1}}}
\newcommand{\AnnotationTok}[1]{\textcolor[rgb]{0.38,0.63,0.69}{\textbf{\textit{#1}}}}
\newcommand{\AttributeTok}[1]{\textcolor[rgb]{0.49,0.56,0.16}{#1}}
\newcommand{\BaseNTok}[1]{\textcolor[rgb]{0.25,0.63,0.44}{#1}}
\newcommand{\BuiltInTok}[1]{\textcolor[rgb]{0.00,0.50,0.00}{#1}}
\newcommand{\CharTok}[1]{\textcolor[rgb]{0.25,0.44,0.63}{#1}}
\newcommand{\CommentTok}[1]{\textcolor[rgb]{0.38,0.63,0.69}{\textit{#1}}}
\newcommand{\CommentVarTok}[1]{\textcolor[rgb]{0.38,0.63,0.69}{\textbf{\textit{#1}}}}
\newcommand{\ConstantTok}[1]{\textcolor[rgb]{0.53,0.00,0.00}{#1}}
\newcommand{\ControlFlowTok}[1]{\textcolor[rgb]{0.00,0.44,0.13}{\textbf{#1}}}
\newcommand{\DataTypeTok}[1]{\textcolor[rgb]{0.56,0.13,0.00}{#1}}
\newcommand{\DecValTok}[1]{\textcolor[rgb]{0.25,0.63,0.44}{#1}}
\newcommand{\DocumentationTok}[1]{\textcolor[rgb]{0.73,0.13,0.13}{\textit{#1}}}
\newcommand{\ErrorTok}[1]{\textcolor[rgb]{1.00,0.00,0.00}{\textbf{#1}}}
\newcommand{\ExtensionTok}[1]{#1}
\newcommand{\FloatTok}[1]{\textcolor[rgb]{0.25,0.63,0.44}{#1}}
\newcommand{\FunctionTok}[1]{\textcolor[rgb]{0.02,0.16,0.49}{#1}}
\newcommand{\ImportTok}[1]{\textcolor[rgb]{0.00,0.50,0.00}{\textbf{#1}}}
\newcommand{\InformationTok}[1]{\textcolor[rgb]{0.38,0.63,0.69}{\textbf{\textit{#1}}}}
\newcommand{\KeywordTok}[1]{\textcolor[rgb]{0.00,0.44,0.13}{\textbf{#1}}}
\newcommand{\NormalTok}[1]{#1}
\newcommand{\OperatorTok}[1]{\textcolor[rgb]{0.40,0.40,0.40}{#1}}
\newcommand{\OtherTok}[1]{\textcolor[rgb]{0.00,0.44,0.13}{#1}}
\newcommand{\PreprocessorTok}[1]{\textcolor[rgb]{0.74,0.48,0.00}{#1}}
\newcommand{\RegionMarkerTok}[1]{#1}
\newcommand{\SpecialCharTok}[1]{\textcolor[rgb]{0.25,0.44,0.63}{#1}}
\newcommand{\SpecialStringTok}[1]{\textcolor[rgb]{0.73,0.40,0.53}{#1}}
\newcommand{\StringTok}[1]{\textcolor[rgb]{0.25,0.44,0.63}{#1}}
\newcommand{\VariableTok}[1]{\textcolor[rgb]{0.10,0.09,0.49}{#1}}
\newcommand{\VerbatimStringTok}[1]{\textcolor[rgb]{0.25,0.44,0.63}{#1}}
\newcommand{\WarningTok}[1]{\textcolor[rgb]{0.38,0.63,0.69}{\textbf{\textit{#1}}}}
\setlength{\emergencystretch}{3em} % prevent overfull lines
\providecommand{\tightlist}{%
  \setlength{\itemsep}{0pt}\setlength{\parskip}{0pt}}
\usepackage{bookmark}
\IfFileExists{xurl.sty}{\usepackage{xurl}}{} % add URL line breaks if available
\urlstyle{same}
\hypersetup{
  hidelinks,
  pdfcreator={LaTeX via pandoc}}

\author{}
\date{}

\begin{document}

\section{Anonymous Equal-Amplitude Composite-Wave Model: Basic Axiom
System
v3}\label{anonymous-equal-amplitude-composite-wave-model-basic-axiom-system-v3}

\textbf{Date:} 2026-07-12\\
\textbf{Author:} Noriaki Kihara\\
\textbf{Status:} Definition paper\\
\textbf{Version DOI:} 10.5281/zenodo.21315736\\
\textbf{Concept DOI:} 10.5281/zenodo.21315735

\begin{center}\rule{0.5\linewidth}{0.5pt}\end{center}

\section{1. First Principles}\label{first-principles}

\subsection{Axiom 0: Component
Anonymity}\label{axiom-0-component-anonymity}

Classification: Axiom

Basic components are not assigned individual names, privileged axes,
privileged signs, or privileged types.

The index \texttt{n} is a convenient label and is not an intrinsic name
of the component.

Forbidden:

\begin{Shaded}
\begin{Highlighting}[]
\NormalTok{making only one specific component the negative{-}sign axis}
\NormalTok{making only one specific component the time axis}
\NormalTok{making only one specific component real{-}valued}
\NormalTok{making only one specific component complex{-}valued}
\NormalTok{giving an intrinsic name to only one specific component}
\end{Highlighting}
\end{Shaded}

\begin{center}\rule{0.5\linewidth}{0.5pt}\end{center}

\subsection{Axiom 0+: Formulation
Anonymity}\label{axiom-0-formulation-anonymity}

Classification: Axiom

At the first-principle layer, theory names, formulation names, and
readout-logic names are not assigned intrinsic names.

Equations, projections, or readout rules at the first-principle layer
must not be changed on the basis of physical names, object names,
interaction names, or existing theory classifications.

Forbidden:

\begin{Shaded}
\begin{Highlighting}[]
\NormalTok{adopting a circular{-}motion equation because it is called gravity}
\NormalTok{adopting a hyperbolic equation because it is called Coulomb force}
\NormalTok{erasing signs because it is called mass}
\NormalTok{preserving signs because it is called charge}
\NormalTok{moving a component to a time axis because it is called energy}
\NormalTok{using a different closure equation because the target is a two{-}body system}
\NormalTok{using a different closure equation because the target is a three{-}body system}
\NormalTok{using a special readout rule because the target is an ABC system}
\NormalTok{choosing a right{-}hand{-}side representation to obtain a desired physical name}
\end{Highlighting}
\end{Shaded}

Allowed:

\begin{Shaded}
\begin{Highlighting}[]
\NormalTok{different readout operations applied to the same closure condition}
\NormalTok{different symmetry breaking applied to the same closure condition}
\NormalTok{different gauge stability applied to the same closure condition}
\NormalTok{different sign preservation applied to the same closure condition}
\NormalTok{different sign erasure applied to the same closure condition}
\NormalTok{different phase{-}branch selection applied to the same closure condition}
\NormalTok{different closure target sets applied to the same closure condition}
\end{Highlighting}
\end{Shaded}

Order:

\begin{Shaded}
\begin{Highlighting}[]
\NormalTok{place the same all{-}positive{-}sign zero closure}
\NormalTok{define the readout operation first}
\NormalTok{define symmetry breaking first}
\NormalTok{define sign preservation or sign erasure first}
\NormalTok{confirm gauge stability}
\NormalTok{assign working names after readout}
\end{Highlighting}
\end{Shaded}

Target set:

\begin{Shaded}
\begin{Highlighting}[]
\NormalTok{Q(P)=\textbackslash{}sum\_n p\_n\^{}2}
\end{Highlighting}
\end{Shaded}

\begin{Shaded}
\begin{Highlighting}[]
\NormalTok{Q(A)=0}
\end{Highlighting}
\end{Shaded}

\begin{Shaded}
\begin{Highlighting}[]
\NormalTok{Q(A+B)=0}
\end{Highlighting}
\end{Shaded}

\begin{Shaded}
\begin{Highlighting}[]
\NormalTok{Q(A+B+C)=0}
\end{Highlighting}
\end{Shaded}

Differences in target set are allowed.

Changing the definition of \texttt{Q} is not allowed.

\begin{center}\rule{0.5\linewidth}{0.5pt}\end{center}

\subsection{Axiom 0.5: Dimension-Wise Central Projection
Readout}\label{axiom-0.5-dimension-wise-central-projection-readout}

Classification: Auxiliary geometric axiom

No curvature correction is introduced for readout along a single axis.

Curvature-correction candidates are introduced only when two or more
independent directions form an area, volume, closed path, or transport
mismatch.

\begin{center}\rule{0.5\linewidth}{0.5pt}\end{center}

\subsection{Axiom 1: All-Positive-Sign Zero
Closure}\label{axiom-1-all-positive-sign-zero-closure}

Classification: Axiom

\begin{Shaded}
\begin{Highlighting}[]
\NormalTok{\textbackslash{}sum\_\{n=1\}\^{}\{N\}x\_n\^{}2=0}
\end{Highlighting}
\end{Shaded}

Expanded form:

\begin{Shaded}
\begin{Highlighting}[]
\NormalTok{x\_1\^{}2+x\_2\^{}2+\textbackslash{}cdots+x\_N\^{}2=0}
\end{Highlighting}
\end{Shaded}

Coefficient signs:

\begin{Shaded}
\begin{Highlighting}[]
\NormalTok{all terms +}
\end{Highlighting}
\end{Shaded}

Not adopted:

\begin{Shaded}
\begin{Highlighting}[]
\NormalTok{\textbackslash{}sum\_n |x\_n|\^{}2=0}
\end{Highlighting}
\end{Shaded}

Term used for closure:

\begin{Shaded}
\begin{Highlighting}[]
\NormalTok{x\_n\^{}2}
\end{Highlighting}
\end{Shaded}

Term not used for closure:

\begin{Shaded}
\begin{Highlighting}[]
\NormalTok{x\_n\textbackslash{}bar\{x\}\_n}
\end{Highlighting}
\end{Shaded}

Applicable states:

\begin{Shaded}
\begin{Highlighting}[]
\NormalTok{stationary closed state}
\NormalTok{stable closed state}
\end{Highlighting}
\end{Shaded}

Quasi-stationary states:

\begin{Shaded}
\begin{Highlighting}[]
\NormalTok{immediately after external influence}
\NormalTok{closure recovery process}
\NormalTok{metastable state}
\end{Highlighting}
\end{Shaded}

Temporary nonzero closure residual in a quasi-stationary state:

\begin{Shaded}
\begin{Highlighting}[]
\NormalTok{\textbackslash{}sum\_n x\_n\^{}2\textbackslash{}ne0}
\end{Highlighting}
\end{Shaded}

After stationarization:

\begin{Shaded}
\begin{Highlighting}[]
\NormalTok{\textbackslash{}sum\_n x\_n\^{}2=0}
\end{Highlighting}
\end{Shaded}

\begin{center}\rule{0.5\linewidth}{0.5pt}\end{center}

\subsection{Axiom 2: Nontrivial
Existence}\label{axiom-2-nontrivial-existence}

Classification: Axiom

\begin{Shaded}
\begin{Highlighting}[]
\NormalTok{\textbackslash{}exists n,\textbackslash{}quad x\_n\textbackslash{}neq0}
\end{Highlighting}
\end{Shaded}

\begin{center}\rule{0.5\linewidth}{0.5pt}\end{center}

\section{2. Consequences from the First
Principles}\label{consequences-from-the-first-principles}

\subsection{Consequence 1: Insufficiency of Real Positive-Definite
Components}\label{consequence-1-insufficiency-of-real-positive-definite-components}

Classification: Derived consequence

\begin{Shaded}
\begin{Highlighting}[]
\NormalTok{x\_n\textbackslash{}in\textbackslash{}mathbb\{R\}}
\NormalTok{\textbackslash{}Rightarrow}
\NormalTok{x\_n\^{}2\textbackslash{}ge0}
\end{Highlighting}
\end{Shaded}

\begin{Shaded}
\begin{Highlighting}[]
\NormalTok{\textbackslash{}sum\_n x\_n\^{}2=0}
\NormalTok{\textbackslash{}Rightarrow}
\NormalTok{x\_n=0\textbackslash{}quad\textbackslash{}forall n}
\end{Highlighting}
\end{Shaded}

\begin{center}\rule{0.5\linewidth}{0.5pt}\end{center}

\subsection{Consequence 2: Internal Generation of Sign
Reversal}\label{consequence-2-internal-generation-of-sign-reversal}

Classification: Derived consequence

\begin{Shaded}
\begin{Highlighting}[]
\NormalTok{x\_1=A,\textbackslash{}qquad x\_2=iA}
\end{Highlighting}
\end{Shaded}

\begin{Shaded}
\begin{Highlighting}[]
\NormalTok{x\_1\^{}2+x\_2\^{}2=A\^{}2+(iA)\^{}2=A\^{}2{-}A\^{}2=0}
\end{Highlighting}
\end{Shaded}

\begin{Shaded}
\begin{Highlighting}[]
\NormalTok{i\^{}2={-}1}
\end{Highlighting}
\end{Shaded}

\begin{center}\rule{0.5\linewidth}{0.5pt}\end{center}

\subsection{Consequence 3: Necessity of Phase
Algebra}\label{consequence-3-necessity-of-phase-algebra}

Classification: Derived consequence

\begin{Shaded}
\begin{Highlighting}[]
\NormalTok{x\_n=r\_n e\^{}\{i\textbackslash{}phi\_n\}}
\end{Highlighting}
\end{Shaded}

\begin{Shaded}
\begin{Highlighting}[]
\NormalTok{x\_n\^{}2=r\_n\^{}2e\^{}\{i2\textbackslash{}phi\_n\}}
\end{Highlighting}
\end{Shaded}

Closure condition:

\begin{Shaded}
\begin{Highlighting}[]
\NormalTok{\textbackslash{}sum\_\{n=1\}\^{}\{N\}r\_n\^{}2e\^{}\{i2\textbackslash{}phi\_n\}=0}
\end{Highlighting}
\end{Shaded}

Equal-amplitude condition:

\begin{Shaded}
\begin{Highlighting}[]
\NormalTok{r\_n=r}
\end{Highlighting}
\end{Shaded}

\begin{Shaded}
\begin{Highlighting}[]
\NormalTok{\textbackslash{}sum\_\{n=1\}\^{}\{N\}e\^{}\{i2\textbackslash{}phi\_n\}=0}
\end{Highlighting}
\end{Shaded}

\begin{center}\rule{0.5\linewidth}{0.5pt}\end{center}

\subsection{Consequence 4: Type
Anonymity}\label{consequence-4-type-anonymity}

Classification: Derived consequence

\begin{Shaded}
\begin{Highlighting}[]
\NormalTok{x\_n\textbackslash{}in\textbackslash{}mathbb\{C\}\textbackslash{}quad\textbackslash{}forall n}
\end{Highlighting}
\end{Shaded}

or

\begin{Shaded}
\begin{Highlighting}[]
\NormalTok{all components are assigned the same phase{-}algebra type}
\end{Highlighting}
\end{Shaded}

\begin{center}\rule{0.5\linewidth}{0.5pt}\end{center}

\subsection{Consequence 5: Ninety-Degree Phase
Difference}\label{consequence-5-ninety-degree-phase-difference}

Classification: Derived consequence

\begin{Shaded}
\begin{Highlighting}[]
\NormalTok{A\^{}2+(iA)\^{}2=0}
\end{Highlighting}
\end{Shaded}

\begin{center}\rule{0.5\linewidth}{0.5pt}\end{center}

\section{3. Anonymous Equal-Amplitude Composite
Wave}\label{anonymous-equal-amplitude-composite-wave}

\subsection{Axiom 3: Anonymous Equal-Amplitude Composite
Wave}\label{axiom-3-anonymous-equal-amplitude-composite-wave}

Classification: Axiom

\begin{Shaded}
\begin{Highlighting}[]
\NormalTok{\textbackslash{}Psi\_N=\textbackslash{}frac\{A\_0\}\{N\}\textbackslash{}sum\_\{k=1\}\^{}\{N\}e\^{}\{i\textbackslash{}theta\_k\}}
\end{Highlighting}
\end{Shaded}

\begin{Shaded}
\begin{Highlighting}[]
\NormalTok{|\textbackslash{}psi\_k|=\textbackslash{}frac\{A\_0\}\{N\}}
\end{Highlighting}
\end{Shaded}

When all phases are aligned:

\begin{Shaded}
\begin{Highlighting}[]
\NormalTok{\textbackslash{}max|\textbackslash{}Psi\_N|=A\_0}
\end{Highlighting}
\end{Shaded}

\begin{center}\rule{0.5\linewidth}{0.5pt}\end{center}

\subsection{Axiom 4: Unreadability of Individual
Phase}\label{axiom-4-unreadability-of-individual-phase}

Classification: Axiom

\begin{Shaded}
\begin{Highlighting}[]
\NormalTok{\textbackslash{}theta\_k\textbackslash{}rightarrow\textbackslash{}theta\_k+\textbackslash{}alpha}
\end{Highlighting}
\end{Shaded}

\begin{center}\rule{0.5\linewidth}{0.5pt}\end{center}

\subsection{Axiom 5: Phase-Difference
Readout}\label{axiom-5-phase-difference-readout}

Classification: Axiom

\begin{Shaded}
\begin{Highlighting}[]
\NormalTok{\textbackslash{}Delta\_\{ij\}=\textbackslash{}theta\_i{-}\textbackslash{}theta\_j}
\end{Highlighting}
\end{Shaded}

Number of readable pair channels:

\begin{Shaded}
\begin{Highlighting}[]
\NormalTok{\textbackslash{}binom\{N\}\{2\}=\textbackslash{}frac\{N(N{-}1)\}\{2\}}
\end{Highlighting}
\end{Shaded}

\begin{center}\rule{0.5\linewidth}{0.5pt}\end{center}

\subsection{Axiom 6: System
Identification}\label{axiom-6-system-identification}

Classification: Axiom

\begin{Shaded}
\begin{Highlighting}[]
\NormalTok{(N,A\_0,\textbackslash{}\{\textbackslash{}Delta\_\{ij\}\textbackslash{}\})}
\end{Highlighting}
\end{Shaded}

\begin{center}\rule{0.5\linewidth}{0.5pt}\end{center}

\subsection{Axiom 7: Composite Phase
Projection}\label{axiom-7-composite-phase-projection}

Classification: Axiom

\begin{Shaded}
\begin{Highlighting}[]
\NormalTok{S\_N=\textbackslash{}sum\_k e\^{}\{i\textbackslash{}theta\_k\}}
\end{Highlighting}
\end{Shaded}

\begin{Shaded}
\begin{Highlighting}[]
\NormalTok{\textbackslash{}Psi\_N=\textbackslash{}frac\{A\_0\}\{N\}S\_N}
\end{Highlighting}
\end{Shaded}

\begin{Shaded}
\begin{Highlighting}[]
\NormalTok{\textbackslash{}Theta\_N=\textbackslash{}arg\textbackslash{}Psi\_N}
\end{Highlighting}
\end{Shaded}

\begin{center}\rule{0.5\linewidth}{0.5pt}\end{center}

\subsection{Axiom 8: No Requirement of Inversion
Symmetry}\label{axiom-8-no-requirement-of-inversion-symmetry}

Classification: Axiom

Inversion symmetry is not required.

\begin{center}\rule{0.5\linewidth}{0.5pt}\end{center}

\section{4. External Readout}\label{external-readout}

\subsection{Axiom 9: External-Axis
Response}\label{axiom-9-external-axis-response}

Classification: Axiom candidate

\begin{Shaded}
\begin{Highlighting}[]
\NormalTok{M\_\textbackslash{}alpha=e\^{}\{i\textbackslash{}alpha\}}
\end{Highlighting}
\end{Shaded}

\begin{Shaded}
\begin{Highlighting}[]
\NormalTok{Y(\textbackslash{}alpha)=\textbackslash{}frac\{A\_0\}\{N\}\textbackslash{}sum\_\{k=1\}\^{}\{N\}e\^{}\{i(\textbackslash{}theta\_k{-}\textbackslash{}alpha)\}}
\NormalTok{=e\^{}\{{-}i\textbackslash{}alpha\}\textbackslash{}Psi\_N}
\end{Highlighting}
\end{Shaded}

Continuous response:

\begin{Shaded}
\begin{Highlighting}[]
\NormalTok{S\_N(\textbackslash{}alpha)=\textbackslash{}operatorname\{Re\}\textbackslash{}left[\textbackslash{}frac\{1\}\{N\}\textbackslash{}sum\_\{k=1\}\^{}\{N\}e\^{}\{i(\textbackslash{}theta\_k{-}\textbackslash{}alpha)\}\textbackslash{}right]}
\end{Highlighting}
\end{Shaded}

Binary response:

\begin{Shaded}
\begin{Highlighting}[]
\NormalTok{S\_N(\textbackslash{}alpha)=\textbackslash{}operatorname\{sgn\}\textbackslash{}operatorname\{Re\}\textbackslash{}left[\textbackslash{}frac\{1\}\{N\}\textbackslash{}sum\_\{k=1\}\^{}\{N\}e\^{}\{i(\textbackslash{}theta\_k{-}\textbackslash{}alpha)\}\textbackslash{}right]}
\end{Highlighting}
\end{Shaded}

\begin{center}\rule{0.5\linewidth}{0.5pt}\end{center}

\subsection{Axiom 10: External Readout
Phase}\label{axiom-10-external-readout-phase}

Classification: Axiom candidate

\begin{Shaded}
\begin{Highlighting}[]
\NormalTok{\textbackslash{}Delta\textbackslash{}phi\_\textbackslash{}mu=\textbackslash{}phi\_\textbackslash{}mu\^{}\{(\textbackslash{}Psi)\}{-}\textbackslash{}phi\_\textbackslash{}mu\^{}\{(M)\}}
\end{Highlighting}
\end{Shaded}

\begin{Shaded}
\begin{Highlighting}[]
\NormalTok{\textbackslash{}mu=x,y,z,t}
\end{Highlighting}
\end{Shaded}

\begin{center}\rule{0.5\linewidth}{0.5pt}\end{center}

\subsection{Axiom 11: Position-Phase
Readout}\label{axiom-11-position-phase-readout}

Classification: Axiom candidate

\begin{Shaded}
\begin{Highlighting}[]
\NormalTok{x,y,z}
\end{Highlighting}
\end{Shaded}

\begin{center}\rule{0.5\linewidth}{0.5pt}\end{center}

\subsection{Axiom 12: Time-Phase
Readout}\label{axiom-12-time-phase-readout}

Classification: Axiom candidate

\begin{Shaded}
\begin{Highlighting}[]
\NormalTok{\textbackslash{}phi\_t\textbackslash{}equiv\textbackslash{}omega t\_\{\textbackslash{}mathrm\{read\}\}\textbackslash{}pmod\{2\textbackslash{}pi\}}
\end{Highlighting}
\end{Shaded}

\begin{center}\rule{0.5\linewidth}{0.5pt}\end{center}

\subsection{Axiom 13: Covering Phase}\label{axiom-13-covering-phase}

Classification: Axiom candidate

\begin{Shaded}
\begin{Highlighting}[]
\NormalTok{2pi{-}periodic readout}
\NormalTok{4pi{-}periodic readout}
\NormalTok{winding{-}number{-}difference readout}
\end{Highlighting}
\end{Shaded}

\begin{center}\rule{0.5\linewidth}{0.5pt}\end{center}

\section{5. Observational Phase
Discreteness}\label{observational-phase-discreteness}

\subsection{Axiom 14: Observational Phase
Discreteness}\label{axiom-14-observational-phase-discreteness}

Classification: Axiom candidate

\begin{Shaded}
\begin{Highlighting}[]
\NormalTok{\textbackslash{}sum\_\{\textbackslash{}mathrm\{cycle\}\}\textbackslash{}Delta\textbackslash{}phi=2\textbackslash{}pi m,\textbackslash{}qquad m\textbackslash{}in\textbackslash{}mathbb\{Z\}}
\end{Highlighting}
\end{Shaded}

\begin{Shaded}
\begin{Highlighting}[]
\NormalTok{W=\textbackslash{}frac\{1\}\{2\textbackslash{}pi\}\textbackslash{}oint d\textbackslash{}phi,\textbackslash{}qquad W\textbackslash{}in\textbackslash{}mathbb\{Z\}}
\end{Highlighting}
\end{Shaded}

Readout-cell form:

\begin{Shaded}
\begin{Highlighting}[]
\NormalTok{\textbackslash{}phi\_\{\textbackslash{}mathrm\{obs\}\}=m\textbackslash{}epsilon\_\textbackslash{}phi,\textbackslash{}qquad m\textbackslash{}in\textbackslash{}mathbb\{Z\}}
\end{Highlighting}
\end{Shaded}

Phase-area candidate:

\begin{Shaded}
\begin{Highlighting}[]
\NormalTok{Q\_\textbackslash{}phi=\textbackslash{}sum\_\{i\textless{}j\}\textbackslash{}sin\^{}2\textbackslash{}frac\{\textbackslash{}Delta\_\{ij\}\}\{2\}}
\end{Highlighting}
\end{Shaded}

\begin{center}\rule{0.5\linewidth}{0.5pt}\end{center}

\section{6. Waveform and Localization}\label{waveform-and-localization}

\subsection{Working Axiom A: State
Waveform}\label{working-axiom-a-state-waveform}

Classification: Working hypothesis

\begin{Shaded}
\begin{Highlighting}[]
\NormalTok{A(\textbackslash{}theta)=\textbackslash{}sum\_n A\_n\textbackslash{}cos(n\textbackslash{}theta+\textbackslash{}phi\_n)}
\end{Highlighting}
\end{Shaded}

\begin{Shaded}
\begin{Highlighting}[]
\NormalTok{Z(\textbackslash{}theta)=\textbackslash{}sum\_n r\_n e\^{}\{i(n\textbackslash{}theta+\textbackslash{}phi\_n)\}}
\end{Highlighting}
\end{Shaded}

\begin{center}\rule{0.5\linewidth}{0.5pt}\end{center}

\subsection{Definition A1: Normalized Equal-Amplitude Odd-Harmonic
Localized
Wave}\label{definition-a1-normalized-equal-amplitude-odd-harmonic-localized-wave}

Classification: Definition

\begin{Shaded}
\begin{Highlighting}[]
\NormalTok{N\_h=2K{-}1}
\end{Highlighting}
\end{Shaded}

\begin{Shaded}
\begin{Highlighting}[]
\NormalTok{S\_\{N\_h\}(u)=\textbackslash{}frac\{1\}\{K\}\textbackslash{}sum\_\{m=0\}\^{}\{K{-}1\}\textbackslash{}cos((2m+1)u)}
\end{Highlighting}
\end{Shaded}

\begin{Shaded}
\begin{Highlighting}[]
\NormalTok{Z\_\{N\_h\}(u)=\textbackslash{}frac\{1\}\{K\}\textbackslash{}sum\_\{m=0\}\^{}\{K{-}1\}e\^{}\{i(2m+1)u\}}
\end{Highlighting}
\end{Shaded}

\begin{Shaded}
\begin{Highlighting}[]
\NormalTok{\textbackslash{}Psi\_\{N\_h\}(u)=A S\_\{N\_h\}(u)}
\end{Highlighting}
\end{Shaded}

\begin{Shaded}
\begin{Highlighting}[]
\NormalTok{\textbackslash{}Psi\_\{N\_h\}(u)=A Z\_\{N\_h\}(u)}
\end{Highlighting}
\end{Shaded}

\begin{Shaded}
\begin{Highlighting}[]
\NormalTok{\textbackslash{}frac\{A\}\{K\}}
\end{Highlighting}
\end{Shaded}

\begin{center}\rule{0.5\linewidth}{0.5pt}\end{center}

\subsection{Theorem A1: Representative-Amplitude
Preservation}\label{theorem-a1-representative-amplitude-preservation}

Classification: Derived theorem

\begin{Shaded}
\begin{Highlighting}[]
\NormalTok{S\_\{N\_h\}(0)=1}
\end{Highlighting}
\end{Shaded}

\begin{Shaded}
\begin{Highlighting}[]
\NormalTok{\textbackslash{}Psi\_\{N\_h\}(0)=A}
\end{Highlighting}
\end{Shaded}

\begin{center}\rule{0.5\linewidth}{0.5pt}\end{center}

\subsection{Theorem A2: Highest Odd-Harmonic Order and Localization
Width}\label{theorem-a2-highest-odd-harmonic-order-and-localization-width}

Classification: Derived theorem

\begin{Shaded}
\begin{Highlighting}[]
\NormalTok{K=\textbackslash{}frac\{N\_h+1\}\{2\}}
\end{Highlighting}
\end{Shaded}

\begin{Shaded}
\begin{Highlighting}[]
\NormalTok{\textbackslash{}lambda\_\{N\_h\}=\textbackslash{}frac\{\textbackslash{}lambda\_0\}\{N\_h\}}
\end{Highlighting}
\end{Shaded}

\begin{Shaded}
\begin{Highlighting}[]
\NormalTok{\textbackslash{}Delta x\_\{N\_h\}\textbackslash{}propto\textbackslash{}frac\{\textbackslash{}lambda\_0\}\{N\_h\}}
\end{Highlighting}
\end{Shaded}

\begin{Shaded}
\begin{Highlighting}[]
\NormalTok{N\_\{h,\textbackslash{}rm req\}\textbackslash{}sim\textbackslash{}frac\{\textbackslash{}lambda\_0\}\{\textbackslash{}Delta x\}}
\end{Highlighting}
\end{Shaded}

\begin{Shaded}
\begin{Highlighting}[]
\NormalTok{N\_\{h,\textbackslash{}rm req\}\^{}\{(t)\}\textbackslash{}sim\textbackslash{}frac\{T\_0\}\{\textbackslash{}Delta t\}}
\end{Highlighting}
\end{Shaded}

\begin{center}\rule{0.5\linewidth}{0.5pt}\end{center}

\subsection{Consequence A1: Identical Construction of Spatial and
Temporal
Localization}\label{consequence-a1-identical-construction-of-spatial-and-temporal-localization}

Classification: Derived consequence

\begin{Shaded}
\begin{Highlighting}[]
\NormalTok{\textbackslash{}Psi(\textbackslash{}chi,\textbackslash{}tau)=A S\_\{N\_\{h,\textbackslash{}chi\}\}(\textbackslash{}chi{-}\textbackslash{}chi\_0)S\_\{N\_\{h,\textbackslash{}tau\}\}(\textbackslash{}tau{-}\textbackslash{}tau\_0)e\^{}\{i\textbackslash{}phi\}}
\end{Highlighting}
\end{Shaded}

\begin{center}\rule{0.5\linewidth}{0.5pt}\end{center}

\subsection{Working Axiom B:
Localization}\label{working-axiom-b-localization}

Classification: Working hypothesis

\begin{Shaded}
\begin{Highlighting}[]
\NormalTok{S\_\{N\_h\}(u)=\textbackslash{}frac\{1\}\{K\}\textbackslash{}sum\_\{m=0\}\^{}\{K{-}1\}\textbackslash{}cos((2m+1)u)}
\end{Highlighting}
\end{Shaded}

\begin{Shaded}
\begin{Highlighting}[]
\NormalTok{Z\_\{N\_h\}(u)=\textbackslash{}frac\{1\}\{K\}\textbackslash{}sum\_\{m=0\}\^{}\{K{-}1\}e\^{}\{i(2m+1)u\}}
\end{Highlighting}
\end{Shaded}

\begin{center}\rule{0.5\linewidth}{0.5pt}\end{center}

\subsection{Working Axiom C: Coupling
Entrance}\label{working-axiom-c-coupling-entrance}

Classification: Working hypothesis

The entrance to interaction is low-order base-phase alignment.

\begin{center}\rule{0.5\linewidth}{0.5pt}\end{center}

\subsection{Working Axiom D: Observation
Map}\label{working-axiom-d-observation-map}

Classification: Working hypothesis

\begin{Shaded}
\begin{Highlighting}[]
\NormalTok{s\_\textbackslash{}alpha=\textbackslash{}mathrm\{Loc\}\textbackslash{}left[\textbackslash{}int A(\textbackslash{}theta)M\_\textbackslash{}alpha(\textbackslash{}theta)d\textbackslash{}theta\textbackslash{}right]}
\end{Highlighting}
\end{Shaded}

\begin{center}\rule{0.5\linewidth}{0.5pt}\end{center}

\subsection{Working Axiom E: Squared Correlation
Strength}\label{working-axiom-e-squared-correlation-strength}

Classification: Working hypothesis

\begin{Shaded}
\begin{Highlighting}[]
\NormalTok{P(\textbackslash{}alpha)=\textbackslash{}frac\{|C\_\textbackslash{}alpha|\^{}2\}\{\textbackslash{}sum\_\textbackslash{}beta |C\_\textbackslash{}beta|\^{}2\}}
\end{Highlighting}
\end{Shaded}

\begin{center}\rule{0.5\linewidth}{0.5pt}\end{center}

\subsection{Working Axiom F: Shared Low-Order Base
Synchronization}\label{working-axiom-f-shared-low-order-base-synchronization}

Classification: Working hypothesis

\begin{Shaded}
\begin{Highlighting}[]
\NormalTok{\textbackslash{}theta\_A{-}\textbackslash{}theta\_B=\textbackslash{}Delta}
\end{Highlighting}
\end{Shaded}

\begin{center}\rule{0.5\linewidth}{0.5pt}\end{center}

\subsection{Working Axiom G: Synchronized Demodulation
Failure}\label{working-axiom-g-synchronized-demodulation-failure}

Classification: Working hypothesis

\begin{Shaded}
\begin{Highlighting}[]
\NormalTok{V(t)=\textbackslash{}left|\textbackslash{}sum\_n w\_n e\^{}\{i\textbackslash{}epsilon\_n(t)\}\textbackslash{}right|}
\end{Highlighting}
\end{Shaded}

\begin{Shaded}
\begin{Highlighting}[]
\NormalTok{V(t)=\textbackslash{}left|\textbackslash{}sum\_n w\_n e\^{}\{{-}\textbackslash{}frac12\textbackslash{}sigma\_n\^{}2(t)\}\textbackslash{}right|}
\end{Highlighting}
\end{Shaded}

\begin{center}\rule{0.5\linewidth}{0.5pt}\end{center}

\subsection{Working Axiom H: Hierarchy}\label{working-axiom-h-hierarchy}

Classification: Working hypothesis

\begin{Shaded}
\begin{Highlighting}[]
\NormalTok{low{-}order base mode}
\NormalTok{+ high{-}order localized mode}
\end{Highlighting}
\end{Shaded}

\begin{center}\rule{0.5\linewidth}{0.5pt}\end{center}

\section{7. Readout, Memory, and Complex
Representation}\label{readout-memory-and-complex-representation}

\subsection{I/Q Readout}\label{iq-readout}

Classification: Theoretical observation

\begin{Shaded}
\begin{Highlighting}[]
\NormalTok{z=a+ib}
\end{Highlighting}
\end{Shaded}

\begin{Shaded}
\begin{Highlighting}[]
\NormalTok{a=r\textbackslash{}cos\textbackslash{}theta,\textbackslash{}qquad b=r\textbackslash{}sin\textbackslash{}theta}
\end{Highlighting}
\end{Shaded}

\begin{Shaded}
\begin{Highlighting}[]
\NormalTok{u=r\textbackslash{}cos\textbackslash{}theta+r\textbackslash{}sin\textbackslash{}theta}
\NormalTok{=\textbackslash{}sqrt2 r\textbackslash{}cos\textbackslash{}left(\textbackslash{}theta{-}\textbackslash{}frac\{\textbackslash{}pi\}\{4\}\textbackslash{}right)}
\end{Highlighting}
\end{Shaded}

\begin{center}\rule{0.5\linewidth}{0.5pt}\end{center}

\subsection{Harmonic-Information
Readout}\label{harmonic-information-readout}

Classification: Theoretical observation

\begin{Shaded}
\begin{Highlighting}[]
\NormalTok{Z(\textbackslash{}theta)=\textbackslash{}sum\_n r\_n e\^{}\{i(n\textbackslash{}theta+\textbackslash{}phi\_n)\}}
\end{Highlighting}
\end{Shaded}

\begin{Shaded}
\begin{Highlighting}[]
\NormalTok{P(Z)=\textbackslash{}operatorname\{Re\}(Z)+\textbackslash{}operatorname\{Im\}(Z)}
\end{Highlighting}
\end{Shaded}

\begin{Shaded}
\begin{Highlighting}[]
\NormalTok{P(Z)=\textbackslash{}sum\_n \textbackslash{}sqrt2 r\_n}
\NormalTok{\textbackslash{}cos\textbackslash{}left(n\textbackslash{}theta+\textbackslash{}phi\_n{-}\textbackslash{}frac\{\textbackslash{}pi\}\{4\}\textbackslash{}right)}
\end{Highlighting}
\end{Shaded}

\begin{center}\rule{0.5\linewidth}{0.5pt}\end{center}

\subsection{Preservation Side and Readout
Side}\label{preservation-side-and-readout-side}

Classification: Theoretical observation

\begin{Shaded}
\begin{Highlighting}[]
\NormalTok{z\textbackslash{}bar z=a\^{}2+b\^{}2}
\end{Highlighting}
\end{Shaded}

\begin{Shaded}
\begin{Highlighting}[]
\NormalTok{z\^{}2=a\^{}2{-}b\^{}2+2iab}
\end{Highlighting}
\end{Shaded}

\begin{center}\rule{0.5\linewidth}{0.5pt}\end{center}

\subsection{Square Registry}\label{square-registry}

Classification: Working hypothesis

\begin{Shaded}
\begin{Highlighting}[]
\NormalTok{\textbackslash{}mathcal\{A\}\_0=\textbackslash{}sum\_k A\_k\^{}2}
\end{Highlighting}
\end{Shaded}

\begin{Shaded}
\begin{Highlighting}[]
\NormalTok{a\_0=\textbackslash{}sum\_k |A\_k|\^{}2=\textbackslash{}sum\_k A\_k\textbackslash{}bar A\_k}
\end{Highlighting}
\end{Shaded}

\begin{center}\rule{0.5\linewidth}{0.5pt}\end{center}

\section{8. Statistical
Classification}\label{statistical-classification}

\subsection{Pair Phase Difference}\label{pair-phase-difference}

Classification: Definition

\begin{Shaded}
\begin{Highlighting}[]
\NormalTok{\textbackslash{}Delta\_\{ij\}=\textbackslash{}theta\_i{-}\textbackslash{}theta\_j}
\end{Highlighting}
\end{Shaded}

\begin{Shaded}
\begin{Highlighting}[]
\NormalTok{\textbackslash{}binom\{N\}\{2\}=\textbackslash{}frac\{N(N{-}1)\}\{2\}}
\end{Highlighting}
\end{Shaded}

Pair composite:

\begin{Shaded}
\begin{Highlighting}[]
\NormalTok{\textbackslash{}psi\_i+\textbackslash{}psi\_j}
\NormalTok{=2A\textbackslash{}cos\textbackslash{}frac\{\textbackslash{}Delta\_\{ij\}\}\{2\}e\^{}\{i(\textbackslash{}theta\_i+\textbackslash{}theta\_j)/2\}}
\end{Highlighting}
\end{Shaded}

Pair interference strength:

\begin{Shaded}
\begin{Highlighting}[]
\NormalTok{I\_\{ij\}=4A\^{}2\textbackslash{}cos\^{}2\textbackslash{}frac\{\textbackslash{}Delta\_\{ij\}\}\{2\}}
\end{Highlighting}
\end{Shaded}

\begin{center}\rule{0.5\linewidth}{0.5pt}\end{center}

\subsection{Overlap Degree}\label{overlap-degree}

Classification: Definition

\begin{Shaded}
\begin{Highlighting}[]
\NormalTok{R\_N=\textbackslash{}left|\textbackslash{}frac\{A\_0\}\{N\}\textbackslash{}sum\_\{k=1\}\^{}\{N\}e\^{}\{i\textbackslash{}theta\_k\}\textbackslash{}right|}
\end{Highlighting}
\end{Shaded}

\begin{Shaded}
\begin{Highlighting}[]
\NormalTok{\textbackslash{}rho\_N=\textbackslash{}frac\{R\_N\}\{A\_0\}}
\NormalTok{=\textbackslash{}frac\{1\}\{N\}\textbackslash{}left|\textbackslash{}sum\_\{k=1\}\^{}\{N\}e\^{}\{i\textbackslash{}theta\_k\}\textbackslash{}right|}
\end{Highlighting}
\end{Shaded}

\begin{Shaded}
\begin{Highlighting}[]
\NormalTok{0\textbackslash{}le\textbackslash{}rho\_N\textbackslash{}le1}
\end{Highlighting}
\end{Shaded}

\begin{center}\rule{0.5\linewidth}{0.5pt}\end{center}

\subsection{Fermionic Degree}\label{fermionic-degree}

Classification: Definition

\begin{Shaded}
\begin{Highlighting}[]
\NormalTok{F\_N=}
\NormalTok{\textbackslash{}frac\{2\}\{N(N{-}1)\}}
\NormalTok{\textbackslash{}sum\_\{i\textless{}j\}\textbackslash{}sin\^{}2\textbackslash{}frac\{\textbackslash{}Delta\_\{ij\}\}\{2\}}
\end{Highlighting}
\end{Shaded}

\begin{Shaded}
\begin{Highlighting}[]
\NormalTok{B\_N=1{-}F\_N}
\end{Highlighting}
\end{Shaded}

\begin{Shaded}
\begin{Highlighting}[]
\NormalTok{B\_N=}
\NormalTok{\textbackslash{}frac\{2\}\{N(N{-}1)\}}
\NormalTok{\textbackslash{}sum\_\{i\textless{}j\}\textbackslash{}cos\^{}2\textbackslash{}frac\{\textbackslash{}Delta\_\{ij\}\}\{2\}}
\end{Highlighting}
\end{Shaded}

\begin{Shaded}
\begin{Highlighting}[]
\NormalTok{F\_N=\textbackslash{}frac\{N\}\{2(N{-}1)\}(1{-}\textbackslash{}rho\_N\^{}2)}
\end{Highlighting}
\end{Shaded}


\end{document}
